% frontmatter/howto.tex
\chapter{如何使用本书}

\section*{先建立一个共同的阅读起点}

建议先读第1—3章。第1章说明什么变化属于本书讨论的自进化，第2章把版本、可更新对象和发布决策落到系统，第3章解释怎样比较更新方案。随后可以按工作问题选择路径，不必把方法谱系当作必须依次采用的技术路线。

第3章的四格试验提供预算分配的具体案例，也保留其历史设计缺口。阅读时应学习它提出的问题和审计方式，不能把已有四格结果直接当作通用因果判据。复制一个表格布局不等于复制了可识别的对照，实际实现中的来源、调用与权限仍需检查。

\section*{按问题选择章节}

\begin{description}
 \item[准备加入经验记忆。] 第4章先检查反馈来源，第5章处理经验抽象、检索与迁移，第6章讨论何时把记忆操作交给可学习策略。读完后应能指出条目来源、适用条件、读取位置和撤销路径，而不只是选择一个存储工具。
 \item[准备优化提示或技能。] 第7章看提示、示例与上下文策略，第8章看可执行技能的构建、准入与复用。两章共同的问题是：保存的产物到底承担什么功能，以及构建成本需要多少后续任务才能分摊。
 \item[准备改变系统流程。] 第9章讨论工作流、角色与结构搜索。应特别核对候选生成、选择评价和最终报告是否分开，并区分并行墙钟与总资源消耗。
 \item[准备使用参数学习。] 第10章比较经验进入参数与留在外部的条件，第11章解释课程、轨迹和交互式学习。附录C补充策略、回报、优势与约束；训练算法的名字不能替代训练数据和目标说明。
 \item[准备让系统修改自身代码。] 第12章区分任务实现、执行器与修改器。重点是可修改范围、固定外层、档案选择和元层贡献的归属，不要仅从代码库被改写就推断所有环节已经自主化。
 \item[准备持续部署。] 第13章处理更新时机与总账，第14章处理漂移、污染和恢复。阅读这两章时，应同时列出可观察信号与恢复对象；检测到异常、隔离版本、恢复文件和恢复行为是不同步骤。
 \item[准备复核当前证据。] 第15章提供从配置到逐次判定的试验审计与迁移设计，第16章用十个研究问题连接已有观察、适用范围和下一项对照。阅读时按轨迹、分析与下一步设计三层记录自己的判断。
\end{description}

如果只需要审查一份外部方案，可以先用六问建立一页方法卡，再读相应机制章，最后回到第3、15、16章核对证据。这条路径适合比较几个候选方案，避免在实现细节里迷失问题。如果准备开展研究，则应先确定估计量、任务覆盖和资源约束，再决定采用哪种机制。

\section*{怎样读一章中的六个环节}

各章尽量从具体失败开始，再提出设计选择，解释机制，比较对照，讨论条件变化，给出选择原则。章首失败不是为了宣判某个方向无用，而是为了让后面的概念有一个可定位的对象。机制说明回答“怎样发生”，对照回答“有什么额外作用”，条件变化回答“换一个设置还剩下什么”。三者需要连接，但不能互相代替。

看到“选择原则”时，先查它依赖哪些条件。一个原则可能是成本算式的直接推论，也可能是根据已知失效提出的工程建议，或是有限数据支持的经验判断。强度不同，使用方式也不同。看到“开放问题”时，则应区分尚未运行、定义尚未冻结和现有数据无法识别；这些缺口需要不同的下一步。

每章习题按“章号-题号”定位，附录G开头列出逐题导航。G1—G18是补充题集，只有列明的对应才可当作相关分析；开放设计题没有唯一标准答案。两套习题都强调判据和可改变答案的条件。参考分析不是唯一措辞，合理替代设计也可能成立，只要它明确所问问题、输入信息、比较对象与限制。运行型题不要求读者把旧数据覆写重来；先做离线审计往往更有助于理解设计。

\section*{三种数字身份，与一种容易遗漏的空白}

\begin{itemize}
 \item \textbf{引用}：来自原论文或其他明确来源的报告值。阅读时保留模型、任务、指标、预算和选择条件；本书引用它不等于本书复现了它。
 \item \textbf{计算}：从已知输入计算出的差值、比例、规划量或教学数值。它的可信度同时取决于输入身份与公式条件；人为教学输入不能被读成真实运行结果。
 \item \textbf{实测}：本书实际运行记录中的观察，能够追到运行标识和工件。实测不自动意味着对照有效，也不自动支持任意任务上的推广。
\end{itemize}

正文在相应数字旁标出“引用”“计算”和“实测”。公式常数、章节编号与日期不都是经验数据，不能把每个阿拉伯数字都理解为效果报告。看到没有被可靠记录的费用、没有执行的恢复或没有运行的实验，应保留其空白；“未知”与“零”有不同意义。

区间与点估计也要一起读。点估计给出本轮观察的方向和大小，区间反映特定抽样与分析假设下的不确定性。区间跨零不证明等价，某个组别的区间与另一组重叠也不直接替代差值检验。第3章和附录A说明这些区别，第16章将它们落实到具体研究问题。

\section*{从一句话找到它的依据}

外部事实先沿正文引文到原论文，再结合来源索引检查具体节、表和版本。本书运行先沿实验标识到配置与轨迹，再检查汇总口径。计算量则查对应脚本与输入。方法卡帮助快速定位机制，不能取代原始来源；摘要帮助快速了解结果，也不能取代失败轨迹。

附录B介绍平台能做什么，附录D给出从文件身份、事件合账到构建快照的离线复核入口，附录F说明成本覆盖。它们可以先被阅读和检查，不需要立即发起付费模型运行。若要重新实验，须另外确认模型、任务与预算范围，并保存新的执行身份；重新运行不会追溯改变旧数据。

读完一条路线后，可以为自己的系统交付一页分析：指出更新对象与实际读取路径，列出共同起点和资源范围，选择一项能够改变采用决策的检验。发现书中问题时，用具体句子、依据和影响范围进入附录G的勘误链路；新的判断也就有了可继续核对的起点。
